\section{Color-frame correspondence}
\label{app:color-correspondence}
The color connection used in Sections~\ref{sec:color} and \ref{sec:energy} can be compared with a connection on the configuration quotient. The comparison must retain the metric, normalization and changing frame on the stated regular domains.

Start with the real six-frame from Section~\ref{sec:color}. Differentiating its orthonormality gives
\AlgebraFrameSkewDisplay
so its connection is skew-symmetric. This is the real-frame identity underlying the compatible color projection.

To compare normalized color descriptions, let $\SpectrumGraphMap$ map the differential graph to the configuration quotient. Write $\SpectrumQuotientMetric$ for its positive Hermitian metric matrix and $\SpectrumGramMatrix$ for the graph Gram matrix. These represent the metrics on the corresponding spaces in the stated common dimension.

Let $\SpectrumNormalizer$ convert the graph coordinates to an orthonormal basis, and write $\SpectrumNormalizedMap$ for the resulting map between normalized coordinates. Their metric relations are
\SpectrumActualStateMapDisplay
The first equality says that pulling back the quotient pairing gives the graph pairing. The matrix $\SpectrumNormalizer$ normalizes the graph coordinates. Multiplication by the positive square root of the quotient metric then gives $\SpectrumNormalizedMap$, whose adjoint product is the identity. It therefore preserves the inner product of those normalized color coordinates.

For diagonal metrics the operations are explicit,
\SpectrumNormalizationExampleDisplay
The normalizer removes the graph's three scale factors before the map is applied. This example illustrates the metric equality without changing its domain or assumptions.

As the source varies, its normalized frame varies too. Writing $\SpectrumQuotientConnection$ for the quotient color connection, with $(q)$ denoting this description, gives
\SpectrumActualConnectionDisplay
The derivative of the normalized map accounts for that frame change. Including it makes curvature transform by conjugation.

The connection entering the full field may also differ from this transported connection. Keeping their difference gives
\SpectrumActualCurvatureDisplay
where $\SpectrumConnectionDifference$. Its covariant derivative contains the interaction with the transported connection,
\SpectrumCorrectionExampleDisplay
and the remaining quadratic term contains the correction's self-product. These are the terms retained when the full curvature enters the energy contraction in Section~\ref{sec:energy}.

